Disaggregated systems require compact, queryable summaries of remote state to resolve the impedance mismatch between stateless compute and network-separated storage. We demonstrated that the system-level value of probabilistic data structures derives from two properties that resolve this mismatch. The first, dynamic mutability, enables state synchronization with evolving datasets through efficient element deletion. The second, composability, enables distributed aggregation without central coordination through operations that are commutative, associative, and idempotent.

The case studies on Cuckoo Filters (\S\ref{sec:membership}) and HyperLogLog (\S\ref{sec:cardinality}) validated this framework. We showed that specific deployment patterns emerge as consequences of these properties: co-partitioning of filter state with data for oracles optimized for mutability, and hierarchical aggregation trees for oracles optimized for composability.

The framework (\S\ref{sec:framework}) establishes that mutability requires element-level information retention while composability requires lossy compression through commutative operations. Theorem~\ref{thm:incompatibility} formalizes this information-theoretic tension: probabilistic oracles with idempotent merge operations cannot support deletion without false negatives unless they maintain reversible element representations\footnote{The theorem assumes oracles without auxiliary metadata structures. Systems may circumvent this constraint by maintaining separate tombstone sets or deletion logs, but such approaches incur space overhead proportional to the number of deleted elements, negating the compactness property that motivates probabilistic structures.}. This structural impossibility explains why examined structures optimize for one property. This structural constraint explains why examined structures optimize for one property. The Cuckoo Filter maintains reversible representations through explicit fingerprints and supports deletion, but its merge operation is not idempotent. The HyperLogLog sketch has an idempotent merge operation through register-wise maximum but discards element identities, precluding deletion.

Structure selection follows from analyzing two workload characteristics: whether the dataset undergoes explicit deletion operations, and whether query results aggregate information from multiple partitions. The intersection of these characteristics determines which property the oracle must satisfy and thus which deployment pattern the system must instantiate.

Workloads requiring explicit element deletion and single-partition queries necessitate structures satisfying dynamic mutability, deployed in co-partitioned architectures where oracle state aligns with data partitioning. Workloads requiring cardinality aggregation across partitions over static or append-with-expiry datasets necessitate structures satisfying composability, deployed in hierarchical reduction architectures. Workloads requiring both properties confront the constraint established in Theorem~\ref{thm:incompatibility}. Systems addressing this workload class maintain separate specialized oracles: a Cuckoo Filter for membership testing with deletion support, and a HyperLogLog sketch for cardinality aggregation. This approach incurs concrete costs: maintaining two structures incurs memory overhead from maintaining both structures, requires coordination between structures when elements are both deleted and aggregated, and complicates failure recovery protocols that must restore consistency across both oracles.

In Theorem~\ref{thm:incompatibility} we have established the tension between idempotent composability and deletion support under strict requirements: perfect idempotence and zero false negatives during deletion. However, relaxations of these requirements may enable structures that approximate both properties within specified error tolerances.

Consider for example allowing approximate deletion, permitting false negative rate $\eta$ for remaining elements, and approximate idempotence where $P(\mathcal{D} \cup \mathcal{D} = \mathcal{D}) \geq 1 - \delta$ for small $\delta > 0$.

Specific open questions then include: can a structure with false negative rate at most $\eta$ during deletion achieve $\delta$-idempotent merge where the probability that merging a sketch with itself produces the original state satisfies $P(\mathcal{D} \cup \mathcal{D} = \mathcal{D}) \geq 1 - \delta$? What space lower bounds exist for structures simultaneously achieving $\eta$-deletion and $\delta$-composability?
